\section{Riesgos materializados y Lecciones aprendidas}

Durante el desarrollo del proyecto se materializó el riesgo asociado a la dependencia de servicios externos, debido a posibles límites de uso, modificaciones en sus condiciones o interrupciones temporales. Esta situación podría afectar la disponibilidad de los servicios de mapas y pagos, así como la continuidad de los flujos principales de la aplicación. Para reducir su impacto, se consideró necesario revisar las condiciones de uso, controlar el consumo de estos servicios y definir alternativas técnicas que permitan mantener una funcionalidad básica en caso de indisponibilidad.

Asimismo, se presentaron dificultades durante la integración de los métodos de pago, principalmente relacionadas con la configuración de las credenciales, la validación de las transacciones y el tratamiento de respuestas exitosas o fallidas. Estas dificultades evidenciaron que la implementación de pagos requiere pruebas específicas de seguridad, compatibilidad y manejo de errores antes de habilitarse para los usuarios finales. Como consecuencia, fue necesario ajustar la planificación y realizar validaciones progresivas mediante prototipos y pruebas controladas.

La principal lección aprendida es la importancia de identificar y documentar desde el inicio las dependencias externas, sus restricciones y sus posibles alternativas. También se comprobó que los flujos críticos deben validarse tempranamente, especialmente los relacionados con pagos, autenticación y disponibilidad de servicios. Finalmente, se reconoce la necesidad de mantener un registro actualizado de riesgos, responsables, medidas de mitigación y resultados obtenidos, de modo que las decisiones tomadas puedan ser revisadas y mejoradas durante las siguientes etapas del proyecto.
